% =====================================================================
%  Annexe B : Aide-memoire
% =====================================================================
\chapter{Aide-mémoire}
\label{ann:aide-memoire}

\chapitreresume{Cette annexe rassemble les commandes, les réglages, les balises et
les méthodes que l'on oublie le plus vite. Gardez-la à portée de main pendant vos
premiers projets.}

\section{Commandes}

\begin{longtable}{L{5.5cm} L{9.7cm}}
\toprule
\textbf{Commande} & \textbf{Effet} \\
\midrule
\endhead
\code{django-admin startproject nom .} & Crée un projet dans le dossier courant. \\
\code{python manage.py startapp app} & Crée une application. \\
\code{python manage.py runserver} & Lance le serveur de développement. \\
\code{python manage.py runserver 8080} & Lance le serveur sur le port 8080. \\
\code{python manage.py makemigrations} & Prépare les migrations. \\
\code{python manage.py migrate} & Applique les migrations. \\
\code{python manage.py showmigrations} & Affiche l'état des migrations. \\
\code{python manage.py sqlmigrate app 0001} & Affiche le SQL d'une migration. \\
\code{python manage.py createsuperuser} & Crée un superutilisateur. \\
\code{python manage.py changepassword user} & Change un mot de passe. \\
\code{python manage.py shell} & Ouvre un shell avec Django configuré. \\
\code{python manage.py test} & Lance les tests. \\
\code{python manage.py collectstatic} & Rassemble les fichiers statiques. \\
\code{python manage.py check} & Vérifie la configuration. \\
\code{python manage.py check --deploy} & Vérifie la sécurité de production. \\
\code{python manage.py dumpdata} & Exporte les données en JSON. \\
\code{python manage.py loaddata fichier.json} & Importe des données. \\
\code{python manage.py flush} & Vide la base, sans supprimer les tables. \\
\code{python -m pip freeze > requirements.txt} & Fige les dépendances. \\
\bottomrule
\end{longtable}

\section{Réglages essentiels}

\begin{longtable}{L{5.6cm} L{9.6cm}}
\toprule
\textbf{Réglage} & \textbf{Rôle} \\
\midrule
\endhead
\code{DEBUG} & Mode développement. \code{False} en production. \\
\code{SECRET\_KEY} & Clé de signature. Secrète, longue, unique. \\
\code{ALLOWED\_HOSTS} & Domaines acceptés. Obligatoire si
\code{DEBUG = False}. \\
\code{INSTALLED\_APPS} & Applications actives. \\
\code{MIDDLEWARE} & Couches de traitement des requêtes. \\
\code{ROOT\_URLCONF} & Module d'URL racine. \\
\code{TEMPLATES} & Moteur et dossiers de gabarits. \\
\code{DATABASES} & Connexion à la base. \\
\code{LANGUAGE\_CODE} & Langue, par exemple \code{fr-fr}. \\
\code{TIME\_ZONE} & Fuseau horaire. \\
\code{STATIC\_URL} & Préfixe des fichiers statiques. \\
\code{STATICFILES\_DIRS} & Dossiers sources des fichiers statiques. \\
\code{STATIC\_ROOT} & Destination de \code{collectstatic}. \\
\code{MEDIA\_URL} & Préfixe des fichiers envoyés. \\
\code{MEDIA\_ROOT} & Dossier de stockage des fichiers envoyés. \\
\code{LOGIN\_URL} & Route de connexion. \\
\code{LOGIN\_REDIRECT\_URL} & Page après connexion. \\
\code{AUTH\_PASSWORD\_VALIDATORS} & Règles de mot de passe. \\
\code{EMAIL\_BACKEND} & Mode d'envoi des courriers. \\
\bottomrule
\end{longtable}

\section{URL et routage}

\begin{codebox}[python]{catalogue/urls.py, modèle}
from django.urls import include, path, re_path
from . import views

app_name = "catalogue"

urlpatterns = [
    path("", views.accueil, name="accueil"),
    path("livres/", views.livre_liste, name="livre_liste"),
    path("livres/<slug:slug>/", views.livre_detail, name="livre_detail"),
    path("auteurs/<int:pk>/", views.auteur_detail, name="auteur_detail"),
    re_path(r"^annee/(?P<annee>[0-9]{4})/$", views.par_annee),
]
\end{codebox}

\begin{longtable}{L{3.8cm} L{11.4cm}}
\toprule
\textbf{Convertisseur} & \textbf{Correspond à} \\
\midrule
\endhead
\code{str} & Toute chaîne sans barre oblique. \\
\code{int} & Entier positif ou nul. \\
\code{slug} & Lettres, chiffres, tirets et soulignés. \\
\code{uuid} & Identifiant UUID. \\
\code{path} & Chaîne pouvant contenir des barres obliques. \\
\bottomrule
\end{longtable}

\begin{codebox}[python]{Construire une URL}
from django.urls import reverse

reverse("catalogue:livre_detail", kwargs={"slug": "le-petit-prince"})
reverse("catalogue:livre_detail", args=["le-petit-prince"])
\end{codebox}

\section{Modèles et champs}

\begin{longtable}{L{4.3cm} L{4.5cm} L{6.1cm}}
\toprule
\textbf{Champ} & \textbf{Argument clé} & \textbf{Usage} \\
\midrule
\endhead
\code{CharField} & \code{max\_length} & Texte court. \\
\code{TextField} & & Texte long. \\
\code{SlugField} & \code{unique} & Identifiant d'URL. \\
\code{EmailField} & & Adresse validée. \\
\code{IntegerField} & \code{default} & Entier. \\
\code{DecimalField} & \code{max\_digits}, \code{decimal\_places} & Montant exact. \\
\code{BooleanField} & \code{default} & Vrai ou faux. \\
\code{DateField} & \code{auto\_now\_add} & Date. \\
\code{DateTimeField} & \code{auto\_now} & Date et heure. \\
\code{ImageField} & \code{upload\_to} & Image, avec Pillow. \\
\code{ForeignKey} & \code{on\_delete}, \code{related\_name} & Un à plusieurs. \\
\code{ManyToManyField} & \code{related\_name} & Plusieurs à plusieurs. \\
\code{OneToOneField} & \code{on\_delete} & Un à un. \\
\bottomrule
\end{longtable}

\begin{longtable}{L{3.3cm} L{11.9cm}}
\toprule
\textbf{\code{on\_delete}} & \textbf{Comportement} \\
\midrule
\endhead
\code{CASCADE} & Supprime les objets liés. \\
\code{PROTECT} & Refuse la suppression s'il existe des objets liés. \\
\code{RESTRICT} & Refuse, avec une sémantique plus récente. \\
\code{SET\_NULL} & Met la clé à \code{NULL} ; exige \code{null=True}. \\
\code{SET\_DEFAULT} & Met la valeur par défaut. \\
\code{DO\_NOTHING} & Ne fait rien : risque d'incohérence. \\
\bottomrule
\end{longtable}

\section{Requêtes ORM}

\begin{codebox}[python]{Les requêtes à connaître}
from django.db.models import Avg, Count, F, Q

Model.objects.all()
Model.objects.filter(champ=valeur)
Model.objects.exclude(champ=valeur)
Model.objects.get(pk=1)
Model.objects.create(champ=valeur)
Model.objects.order_by("-date")
Model.objects.filter(champ__icontains="mot")
Model.objects.filter(date__gte="2024-01-01")
Model.objects.filter(champ__in=[1, 2, 3])
Model.objects.filter(champ__isnull=True)
Model.objects.filter(Q(a=1) | Q(b=2))
Model.objects.select_related("fk")
Model.objects.prefetch_related("m2m")
Model.objects.annotate(total=Count("enfants"))
Model.objects.aggregate(moyenne=Avg("note"))
Model.objects.values("champ")
Model.objects.values_list("champ", flat=True)
Model.objects.count()
Model.objects.exists()
Model.objects.first()
Model.objects.last()
Model.objects.update_or_create(pk=1, defaults={"champ": valeur})
Model.objects.get_or_create(champ=valeur, defaults={...})
instance.save()
instance.delete()
\end{codebox}

\begin{codebox}[python]{Méthodes de QuerySet}
qs.filter(...)
qs.exclude(...)
qs.order_by("champ")
qs.distinct()
qs.values("a", "b")
qs.annotate(...)
qs.aggregate(...)
qs.exists()
qs.count()
qs.first()
qs.last()
qs[:10]
qs.iterator()
\end{codebox}

\section{Balises et filtres de gabarit}

\begin{longtable}{L{4.6cm} L{10.6cm}}
\toprule
\textbf{Balise} & \textbf{Rôle} \\
\midrule
\endhead
\code{\{\% if ... \%\}} & Condition. \\
\code{\{\% for x in liste \%\}} & Boucle, avec \code{\{\% empty \%\}}. \\
\code{\{\% with x=... \%\}} & Alias. \\
\code{\{\% url "nom" arg \%\}} & Construit une URL. \\
\code{\{\% include "f.html" \%\}} & Inclut un gabarit. \\
\code{\{\% extends "base.html" \%\}} & Héritage. \\
\code{\{\% block nom \%\}} & Zone remplaçable. \\
\code{\{\% csrf\_token \%\}} & Jeton de formulaire. \\
\code{\{\% load static \%\}} & Charge les fichiers statiques. \\
\code{\{\% static "css/a.css" \%\}} & URL d'un fichier statique. \\
\code{\{\% load biblio \%\}} & Charge vos balises. \\
\code{\{\% comment \%\}} & Bloc ignoré. \\
\code{\{\% verbatim \%\}} & Affiche les accolades telles quelles. \\
\code{\{\% now "d/m/Y" \%\}} & Date courante. \\
\code{\{\% firstof a b c \%\}} & Première valeur non vide. \\
\code{\{\% cycle "a" "b" \%\}} & Alterne des valeurs. \\
\bottomrule
\end{longtable}

\begin{longtable}{L{4.6cm} L{10.6cm}}
\toprule
\textbf{Filtre} & \textbf{Effet} \\
\midrule
\endhead
\code{length} & Taille d'une liste ou d'une chaîne. \\
\code{upper}, \code{lower} & Casse. \\
\code{capfirst}, \code{title} & Majuscules. \\
\code{truncatechars:80} & Coupe à 80 caractères. \\
\code{truncatewords:15} & Coupe à 15 mots. \\
\code{date:"d/m/Y"} & Formate une date. \\
\code{time:"H:i"} & Formate une heure. \\
\code{floatformat:2} & Deux décimales. \\
\code{default:"..."} & Valeur de remplacement. \\
\code{default\_if\_none:"..."} & Remplace seulement \code{None}. \\
\code{join:", "} & Joint une liste. \\
\code{pluralize} & Ajoute un \og s \fg{} si nécessaire. \\
\code{yesno:"Oui,Non"} & Traduit un booléen. \\
\code{linebreaks} & Transforme les sauts de ligne en paragraphes. \\
\code{escape} & Échappe le HTML. \\
\code{safe} & Désactive l'échappement. À éviter. \\
\code{slugify} & Transforme en slug. \\
\code{striptags} & Supprime les balises HTML. \\
\bottomrule
\end{longtable}

\section{Formulaires}

\begin{codebox}[python]{Form et ModelForm}
from django import forms
from .models import Livre


class ContactForm(forms.Form):
    nom = forms.CharField(max_length=80)
    email = forms.EmailField()
    message = forms.CharField(widget=forms.Textarea)


class LivreForm(forms.ModelForm):
    class Meta:
        model = Livre
        fields = ["titre", "auteur", "annee_publication"]
        widgets = {
            "annee_publication": forms.NumberInput(attrs={"min": 1450}),
        }
\end{codebox}

\begin{codebox}[python]{Traiter un formulaire}
def vue(request):
    if request.method == "POST":
        form = MonForm(request.POST, request.FILES)
        if form.is_valid():
            objet = form.save()
            return redirect(objet)
    else:
        form = MonForm()
    return render(request, "page.html", {"form": form})
\end{codebox}

\begin{longtable}{L{5.2cm} L{10cm}}
\toprule
\textbf{Élément} & \textbf{Rôle} \\
\midrule
\endhead
\code{form.is\_valid()} & Valide et renvoie un booléen. \\
\code{form.cleaned\_data} & Valeurs nettoyées. \\
\code{form.errors} & Erreurs par champ. \\
\code{form.non\_field\_errors} & Erreurs globales. \\
\code{clean\_<champ>()} & Validation d'un champ. \\
\code{clean()} & Validation globale. \\
\code{form.save(commit=False)} & Prépare sans enregistrer. \\
\code{form.save\_m2m()} & Enregistre les relations plusieurs à plusieurs. \\
\code{formset} & Groupe de formulaires identiques. \\
\bottomrule
\end{longtable}

\section{Authentification et permissions}

\begin{codebox}[python]{Fonctions et décorateurs}
from django.contrib.auth import authenticate, login, logout
from django.contrib.auth.decorators import (
    login_required,
    permission_required,
)
from django.contrib.auth.mixins import (
    LoginRequiredMixin,
    PermissionRequiredMixin,
)

user = authenticate(request, username=u, password=p)
login(request, user)
logout(request)

@login_required
def vue(request):
    ...

@permission_required("catalogue.add_livre", raise_exception=True)
def vue2(request):
    ...
\end{codebox}

\begin{longtable}{L{6.0cm} L{9.2cm}}
\toprule
\textbf{Élément} & \textbf{Rôle} \\
\midrule
\endhead
\code{request.user} & Utilisateur courant. \\
\code{request.user.is\_authenticated} & Vrai si connecté. \\
\code{request.user.is\_staff} & Accès à l'administration. \\
\code{request.user.is\_superuser} & Tous les droits. \\
\code{user.has\_perm("app.code")} & Vérifie une permission. \\
\code{perms.app.code} & Vérification dans un gabarit. \\
\code{Group} & Groupe de permissions. \\
\code{user.groups.add(groupe)} & Ajoute un utilisateur à un groupe. \\
\bottomrule
\end{longtable}

\section{Administration}

\begin{codebox}[python]{ModelAdmin minimal}
from django.contrib import admin
from .models import Livre


@admin.register(Livre)
class LivreAdmin(admin.ModelAdmin):
    list_display = ("titre", "auteur", "statut")
    list_filter = ("statut",)
    search_fields = ("titre", "isbn")
    ordering = ("titre",)
    list_select_related = ("auteur",)
    date_hierarchy = "cree_le"
    readonly_fields = ("cree_le", "modifie_le")
\end{codebox}

\section{Codes HTTP}

\begin{longtable}{L{2.4cm} L{12.8cm}}
\toprule
\textbf{Code} & \textbf{Signification} \\
\midrule
\endhead
\code{200} & OK. \\
\code{201} & Créé. \\
\code{204} & Aucun contenu. \\
\code{301} & Déplacé de façon permanente. \\
\code{302} & Redirection temporaire. \\
\code{304} & Non modifié, cache valide. \\
\code{400} & Requête invalide. \\
\code{401} & Authentification requise. \\
\code{403} & Accès interdit. \\
\code{404} & Ressource introuvable. \\
\code{405} & Méthode non autorisée. \\
\code{409} & Conflit. \\
\code{429} & Trop de requêtes. \\
\code{500} & Erreur interne du serveur. \\
\code{502} & Mauvaise passerelle. \\
\code{503} & Service indisponible. \\
\bottomrule
\end{longtable}

\section{Vues basées sur les classes}

\begin{longtable}{L{5.2cm} L{10cm}}
\toprule
\textbf{Classe} & \textbf{Rôle} \\
\midrule
\endhead
\code{TemplateView} & Affiche un gabarit. \\
\code{ListView} & Liste paginée. \\
\code{DetailView} & Fiche d'un objet. \\
\code{CreateView} & Création avec \code{ModelForm}. \\
\code{UpdateView} & Modification. \\
\code{DeleteView} & Suppression avec confirmation. \\
\code{FormView} & Formulaire non lié à un modèle. \\
\code{RedirectView} & Redirection. \\
\bottomrule
\end{longtable}

\begin{codebox}[python]{Exemple de ListView}
from django.views.generic import ListView
from .models import Livre


class LivreListView(ListView):
    model = Livre
    template_name = "catalogue/livres/liste.html"
    context_object_name = "livres"
    paginate_by = 6

    def get_queryset(self):
        return Livre.objects.select_related("auteur")
\end{codebox}
